\noindent\textbf{1. Experimental Setup}

We conducted a comprehensive evaluation of the proposed TSAM method for the Referring Audio-Visual Segmentation (Ref-AVS) task.

\begin{itemize}
    \item \textbf{Datasets:}
    \begin{itemize}
        \item \textbf{Ref-AVS Dataset:} We utilized the Ref-AVS dataset containing 20,000 text expressions and pixel-level annotations across 4,000 10-second videos.
        \item \textbf{Object Categories:} The dataset included audible objects (20 musical instruments, 8 animals, 15 machines, 5 humans) and 3 categories of static, inaudible objects.
        \item \textbf{Splits:}
        \begin{itemize}
            \item \textit{Training Set:} 2,908 videos.
            \item \textit{Validation Set:} 276 videos.
            \item \textit{Test Set:} 818 videos total, subdivided into:
            \begin{itemize}
                \item \textit{Seen:} 292 videos (categories present in training).
                \item \textit{Unseen:} 269 videos (categories not seen during training; tests generalization).
                \item \textit{Null:} 257 videos (text refers to non-existent/not visible objects; empty true masks).
            \end{itemize}
        \end{itemize}
    \end{itemize}
    
    \item \textbf{Evaluation Metrics:}
    \begin{itemize}
        \item \textbf{Standard Metrics:} We employed the Jaccard Index ($\mathcal{J}$) and F-score ($\mathcal{F}$) for the Seen and Unseen subsets.
        \item \textbf{Null Subset Metric:} We employed the metric $\mathcal{S}$, which measures the ratio of the predicted mask area to the background area ($\mathcal{S} = \sqrt{\text{predicted mask area} / \text{background area}}$). A lower $\mathcal{S}$ value indicates better performance (less incorrect segmentation).
    \end{itemize}
    
    \item \textbf{Baselines Compared:}
    \begin{itemize}
        \item \textbf{AVS Methods (Augmented with Text):} AVSBench (PVT-v2 backbone), AVSegFormer (PVT-v2 backbone), GAVS (SAM backbone), SAMA (SAM backbone). \textit{Note: We re-evaluated SAMA by integrating text with audio/visual inputs.}
        \item \textbf{Ref-VOS Methods (Augmented with Audio):} ReferFormer (Video-Swin backbone), R2-VOS (Video-Swin backbone).
        \item \textbf{Ref-AVS Method:} EEMC (Mask2Former backbone; previous state-of-the-art).
    \end{itemize}
    
    \item \textbf{Implementation Details:}
    \begin{itemize}
        \item \textbf{Model Initialization:} We initialized TSAM using the pre-trained ViT-B variant of SAM ($N=12$, embedding dimension $d_{emb}=256$).
        \item \textbf{Architecture Settings:} The temporal branch consisted of $M=4$ blocks. The number of selected audio queries was set to $k=5$.
        \item \textbf{Hardware:} Training was performed on eight AMD GPUs in a distributed setup.
        \item \textbf{Optimizer:} AdamW optimizer was used.
        \item \textbf{Hyperparameters:} Initial learning rate of $1 \cdot 10^{-4}$, batch size of 1.
        \item \textbf{Training Duration:} The model was trained for 15 epochs with periodic evaluations.
        \item \textbf{Model Selection:} We selected the best-performing model on the validation subset for testing.
    \end{itemize}
\end{itemize}

\vspace{1em}
\noindent\textbf{2. Raw Numeric Data}

\vspace{0.5em}
\noindent\textbf{Table 1: Performance comparison on the Ref-AVS dataset}\\
\textit{Note: Baselines marked with \dag{} had text integration added; baselines marked with \ddag{} had audio integration added; * marks our re-implementation of SAMA with text added.}

\begin{center}
\resizebox{\linewidth}{!}{
\begin{tabular}{lllccccc}
\toprule
\textbf{Method} & \textbf{Task} & \textbf{Visual Backbone} & \textbf{Seen J(\%)} & \textbf{Seen F} & \textbf{Unseen J(\%)} & \textbf{Unseen F} & \textbf{Null S($\downarrow$)} \\
\midrule
AVSBench & AVS\dag & PVT-v2 & 23.20 & 0.511 & 32.36 & 0.547 & 0.208 \\
AVSegFormer & AVS\dag & PVT-v2 & 33.47 & 0.470 & 36.05 & 0.501 & 0.171 \\
GAVS & AVS\dag & SAM & 28.93 & 0.498 & 29.82 & 0.497 & 0.190 \\
SAMA & AVS* & SAM & 39.22 & 0.562 & 47.50 & 0.566 & 0.130 \\
ReferFormer & Ref-VOS\ddag & V-Swin & 31.31 & 0.501 & 30.40 & 0.488 & 0.176 \\
R2-VOS & Ref-VOS\ddag & V-Swin & 25.01 & 0.410 & 27.93 & 0.498 & 0.183 \\
EEMC & Ref-AVS & M2F & 34.20 & 0.513 & 49.54 & 0.648 & \textbf{0.007} \\
\textbf{TSAM (Ours)} & Ref-AVS & SAM & \textbf{43.43} & \textbf{0.568} & \textbf{54.58} & \textbf{0.664} & 0.017 \\
\bottomrule
\end{tabular}
}
\end{center}

\vspace{1em}
\noindent\textbf{Table 2: Ablation study on TSAM components and IoU loss}\\
\textit{Legend: TB=Temporal Branch, TMFL=Temporal Modality Fusion Layer, DPM=Dense Prompting Module, SPM=Sparse Prompting Module, CM=Cached Memory, AM=Adapter Module. Mix(S+U) is the average of Seen and Unseen.}

\begin{center}
\resizebox{\linewidth}{!}{
\begin{tabular}{lccccccc}
\toprule
\textbf{Setting} & \textbf{Seen J(\%)} & \textbf{Seen F} & \textbf{Unseen J(\%)} & \textbf{Unseen F} & \textbf{Mix(S+U) J(\%)} & \textbf{Mix(S+U) F} & \textbf{Null S($\downarrow$)} \\
\midrule
(1) TSAM (Full) & \textbf{43.43} & 0.568 & \textbf{54.58} & \textbf{0.664} & \textbf{49.01} & \textbf{0.616} & 0.017 \\
(2) - TB & 33.05 & 0.507 & 50.48 & 0.657 & 41.77 & 0.582 & 0.505 \\
(3) - TMFL & 40.35 & 0.579 & 45.54 & 0.627 & 42.95 & 0.603 & 0.018 \\
(4) - DPM & 42.72 & 0.580 & 49.10 & 0.647 & 45.91 & 0.614 & 0.018 \\
(5) - SPM & 43.04 & 0.580 & 49.75 & 0.652 & 46.40 & \textbf{0.616} & 0.018 \\
(6) - SPM+DPM & 42.60 & 0.602 & 40.58 & 0.604 & 41.59 & 0.603 & 0.018 \\
(7) - CM(a+v) & 42.07 & 0.544 & 49.11 & 0.659 & 45.59 & 0.602 & 0.018 \\
(8) - CM(v) & 42.75 & 0.549 & 51.18 & 0.660 & 46.97 & 0.605 & 0.018 \\
(9) - AM & 43.13 & \textbf{0.600} & 40.79 & 0.617 & 41.96 & 0.609 & 0.017 \\
(10) - $\mathcal{L}_{\mathrm{IoU}}$ & 38.29 & 0.564 & 42.15 & 0.631 & 40.22 & 0.598 & \textbf{0.008} \\
\bottomrule
\end{tabular}
}
\end{center}

\vspace{1em}
\noindent\textbf{Table 3: Effect of audio queries ($k$) and temporal branch depth ($M$)}

\begin{center}
\resizebox{\linewidth}{!}{
\begin{tabular}{lccccccc}
\toprule
\textbf{Variation} & \textbf{Seen J(\%)} & \textbf{Seen F} & \textbf{Unseen J(\%)} & \textbf{Unseen F} & \textbf{Mix(S+U) J(\%)} & \textbf{Mix(S+U) F} & \textbf{Null S($\downarrow$)} \\
\midrule
\textbf{$k=3$} & 43.58 & \textbf{0.579} & 50.43 & 0.655 & 47.01 & \textbf{0.617} & 0.018 \\
\textbf{$k=5$} & 43.43 & 0.568 & \textbf{54.58} & \textbf{0.664} & \textbf{49.01} & 0.616 & \textbf{0.017} \\
\textbf{$k=7$} & 43.24 & 0.573 & 46.95 & 0.630 & 45.10 & 0.601 & 0.018 \\
\textbf{$M=2$} & 42.04 & 0.566 & 53.51 & 0.651 & 47.76 & 0.609 & 0.023 \\
\textbf{$M=4$} & \textbf{43.43} & 0.568 & \textbf{54.58} & \textbf{0.664} & \textbf{49.01} & \textbf{0.616} & \textbf{0.017} \\
\textbf{$M=6$} & 43.27 & \textbf{0.575} & 49.86 & 0.640 & 46.57 & 0.608 & 0.020 \\
\bottomrule
\end{tabular}
}
\end{center}

\vspace{1em}
\noindent\textbf{3. Qualitative Observations}

\vspace{0.5em}
\noindent\textbf{Comparisons with State-of-the-Art:}
\begin{itemize}
    \item \textbf{Backbone Analysis:} We observed that methods utilizing prior segmentation visual backbones (SAM and Mask2Former) generally outperformed those based on PVT-v2 and V-Swin backbones.
    \item \textbf{SAM-Based Baseline Limitations:} Although GAVS and SAMA are SAM-based, they performed worse than EEMC. We noted that SAMA failed to leverage SAM's flexible, promptable nature, and GAVS lacked cohesive multimodal fusion.
    \item \textbf{TSAM Performance:} TSAM achieved the highest performance on both Seen and Unseen test sets. Specifically, TSAM improved over EEMC by 9.23\% in Jaccard Index on the Seen set and 5.04\% on the Unseen set.
    \item \textbf{Null Set Performance:} TSAM fell slightly behind EEMC on the Null test set (S value 0.017 vs 0.007). We attribute this to SAM's inherent limitation of always attempting to produce a segmentation mask even when no target object is present.
\end{itemize}

\vspace{0.5em}
\noindent\textbf{Ablation Observations:}
\begin{itemize}
    \item \textbf{Temporal Branch (TB):} Removing the temporal branch caused the most significant performance degradation (Seen Jaccard dropped from 43.43\% to 33.05\%). This confirmed the critical role of temporal modeling for generalizing across video frames.
    \item \textbf{Prompting Modules (SPM/DPM):} We found that the Sparse Prompting Module (SPM) and Dense Prompting Module (DPM) play complementary roles. Removing both simultaneously resulted in a notable decrease in segmentation performance.
    \item \textbf{Cached Memory (CM):} Omitting cached memory, especially the visual-only memory, degraded performance, highlighting the importance of shared memory for temporal alignment.
    \item \textbf{Adapter Module (AM):} The omission of the adapter module caused a significant performance drop, particularly on the Unseen test set, validating its role in facilitating deep multimodal interaction.
    \item \textbf{Loss Function:} Including the IoU loss ($\mathcal{L}_{\mathrm{IoU}}$) improved performance on Seen and Unseen sets but slightly degraded the Null score.
\end{itemize}

\vspace{0.5em}
\noindent\textbf{Hyperparameter Sensitivity:}
\begin{itemize}
    \item \textbf{Audio Queries ($k$):} We found that $k=5$ yielded optimal results. A lower value ($k=3$) performed well on Seen data but worse on Unseen, while a higher value ($k=7$) likely introduced irrelevant queries, overloading the decoder.
    \item \textbf{Temporal Depth ($M$):} $M=4$ blocks provided the best balance. Shallower setups ($M=2$) lacked sufficient temporal depth, while deeper setups ($M=6$) appeared to introduce excessive complexity that impaired generalization.
\end{itemize}

\vspace{0.5em}
\noindent\textbf{Visual Qualitative Analysis:}
\begin{itemize}
    \item \textbf{Seen Test Set:} We observed that TSAM consistently produced high-quality masks for targeted objects. It successfully segmented inaudible objects when guided by textual cues (e.g., ``The object behind the sounding women''), demonstrating effective processing of complex multimodal instructions.
    \item \textbf{Unseen Test Set:} TSAM demonstrated a remarkable capacity to generalize to novel objects. It effectively aligned audio-visual and textual inputs in novel scenes (e.g., segmenting a ``truck moving'' or ``tuba being played'' that were not in training categories).
    \item \textbf{Generalization:} The visual results underscored TSAM's ability to preserve SAM's pre-trained knowledge while using the temporal branch to understand dynamic scenes.
\end{itemize}